\section{Introduction}

Much top-level RTL and subsystem code does not implement new datapath logic; it connects existing modules, wrappers, channels, buses, and clock/reset infrastructure. A ready/valid edge in a diagram expands into payload, valid, ready, sideband, reset, and sometimes CDC or protocol-adapter wiring. For LLMs this is a poor target: free-form RTL requires long-range consistency over names, directions, widths, protocols, domains, resets, and validation collateral.

The resulting failures are systematic rather than cosmetic. In the benchmark
records, unstructured prompts most often fail by omitting ports, reversing
producer/consumer direction, crossing clock domains without an explicit
adapter, emitting syntactically plausible but semantically rejected wrappers, or
returning JSON that cannot be machine-checked. MICO's causal hypothesis is that
the target representation should expose these obligations as typed graph edges
and contract checks before any RTL is accepted.

MICO takes a narrower position than general RTL generation benchmarks~\cite{verilogeval2023,rtlcoder2023,openllmrtl2025}, domain-adapted chip LLMs~\cite{chipnemo2023}, and recent checked prompt languages~\cite{cppl2026}. The LLM proposes a module-composition graph and optional repair patch; the compiler is the authority. The compiler checks interfaces, roles, widths, protocols, clock domains, adapters, and v0 ready/valid contracts before emitting RTL and traceability artifacts.

This paper makes five contributions:

\begin{enumerate}
  \item a typed module-interface-contract representation that separates LLM proposal from compiler authority;
  \item the MICO language and Rust compiler path for checked graphs, diagnostics, JSON AST repair, SV/SVA emit, and traceability;
  \item ModuleComposeBench, with public-development, held-out, and supplemental realism manifests, positive/negative tasks, case studies, simulation, formal smoke, and QoR collateral;
  \item an authenticated full-matrix evaluation showing that schema-guided JSON AST composition and a bounded compiler-feedback fallback improve tested LLM integration correctness and unsafe rejection over direct RTL, SystemVerilog-interface, and MICO-source prompting;
  \item a reproducible artifact with Docker-first validation, schema gates, release hashes, a Vivado host-tool exception, and reviewer bundle metadata.
\end{enumerate}

The positive LLM claim is deliberately narrow. It covers the tested profiles,
prompts, manifests, schemas, and compiler gates; it does not make the LLM a
source of correctness. The artifact does not replace HDLs or downstream
verification. It attacks the integration layer where engineers already rely on
structure, naming discipline, and review checklists, and where LLMs are most
exposed to textual consistency failures.
